\section{Problem Description}
\label{sec:problem}
% ~0.5 page. Models: hour-only baseline, Phase 1, Phase 3b (final).

Reference stations measure urban air pollution accurately but are expensive,
so cities have few of them. Low-cost metal-oxide sensors could fill the gaps,
but they are noisy, drift, and respond to weather and to gases other than their
target, so they must be calibrated in the field against a reference
\citep{devito2008calibration}. We treat calibration as probabilistic
inference: the true concentration is latent, and the sensors are noisy,
indirect observations of it. A co-located reference analyzer provides ground
truth for training, but its CO, NO$_x$ and NO$_2$ records are missing on
17--18\% of hours (\cref{sec:data}). A Bayesian network suits this setting:
the dependencies among gases, sensors and weather form a directed graph, EM
\citep{dempster1977em} learns from the partially observed references, and
exact inference returns a posterior over each pollution level.

\paragraph{Research questions.} Using only what a deployed device observes
(hour, weather and its own sensors), never a reference reading:
\textbf{Q1 (estimation):} can a discrete BN learned with EM estimate pollution
better than the hour of day alone, and how close does it come to the limit set
by discretization? \textbf{Q2 (shared structure):} the gases come from the
same traffic and rise and fall together; does a latent emission regime capture
that co-movement, and does it improve the estimates?

\paragraph{Approach and findings.} We compare an hour-of-day baseline, a
single-gas network for CO (Phase~1) and our final model (Phase~3b), which
estimates CO, NO$_x$ and NO$_2$ jointly through a latent regime, with sensor
edges chosen by a cross-sensitivity study (\cref{app:cross}). On held-out test
weeks the final model cuts the baseline's error by 37\% for CO, 32\% for
NO$_x$ and 18\% for NO$_2$, closing 87\% and 71\% of the gap to the
discretization floor for CO and NO$_x$ but only 39\% for NO$_2$. The regime
explains most of the gases' co-movement but mainly improves NO$_2$.